\section{Data, code, and AI-use statements}

\paragraph{Data and code availability.}
The complete canonical verification release accompanies this manuscript as \texttt{Erdos506.zip}.  The public repository is \url{https://github.com/DenisUkranian/erdos-506-circles}; DOI and arXiv identifiers will be added after archival deposit.  The release SHA-256 and reproduction instructions are included with the archive.

\paragraph{Generative-AI use.}
Generative-AI systems were used interactively during exploratory proof search, program drafting, debugging, verification orchestration, archival organization and language editing.  They are not authors.  The listed human author(s) reviewed the released material and accept responsibility for every mathematical statement, reference, program and file.  No computational result is accepted solely on the basis of an AI-generated assertion.

\paragraph{Competing interests.}
The author declares no competing interests.

\paragraph{Acknowledgements.}
The author thanks independent readers and maintainers of the Erd\H{o}s Problems resources for preserving the problem and its historical context.  Additional acknowledgements should be inserted here after external verification.

\begin{thebibliography}{99}

\bibitem{erdos1961}
P.~Erd\H{o}s,
\emph{Some unsolved problems},
A Magyar Tudom\'anyos Akad\'emia Matematikai Kutat\'o Int\'ezet\'enek K\"ozlem\'enyei \textbf{6} (1961), 221--254.

\bibitem{elliott1967}
P.~D.~T.~A. Elliott,
\emph{On the number of circles determined by $n$ points},
Acta Math. Acad. Sci. Hungar. \textbf{18} (1967), 181--188,
doi:10.1007/BF02020972.

\bibitem{purdysmith2010}
G.~B. Purdy and J.~W. Smith,
\emph{Lines, circles, planes and spheres},
Discrete Comput. Geom. \textbf{44} (2010), 860--882,
doi:10.1007/s00454-010-9270-3; arXiv:0907.0724.

\bibitem{burr1974}
S.~A. Burr, B. Gr\"unbaum and N.~J.~A. Sloane,
\emph{The orchard problem},
Geom. Dedicata \textbf{2} (1974), 397--424,
doi:10.1007/BF00147569.

\bibitem{kellymoser1958}
L.~M. Kelly and W.~O.~J. Moser,
\emph{On the number of ordinary lines determined by $n$ points},
Canad. J. Math. \textbf{10} (1958), 210--219,
doi:10.4153/CJM-1958-024-6.

\bibitem{csimasawyer1993}
J. Csima and E.~T. Sawyer,
\emph{There exist $6n/13$ ordinary points},
Discrete Comput. Geom. \textbf{9} (1993), 187--202,
doi:10.1007/BF02189318.

\bibitem{langer2003}
A. Langer,
\emph{Logarithmic orbifold Euler numbers of surfaces with applications},
Proc. Lond. Math. Soc. \textbf{86} (2003), 358--396,
doi:10.1112/S0024611502013874.

\bibitem{dezeeuw2018}
F. de Zeeuw,
\emph{Spanned lines and Langer's inequality},
arXiv:1802.08015 (2018).

\bibitem{cuntz2012}
M. Cuntz,
\emph{Simplicial arrangements with up to 27 lines},
Discrete Comput. Geom. \textbf{48} (2012), 682--701,
doi:10.1007/s00454-012-9423-7; arXiv:1108.3000.

\bibitem{shnurnikov2016}
I.~N. Shnurnikov,
\emph{A $t_k$ inequality for arrangements of pseudolines},
Discrete Comput. Geom. \textbf{55} (2016), 284--295,
doi:10.1007/s00454-015-9744-4.

\bibitem{dumnicki2016}
M. Dumnicki, L. Farnik, A. G\l{}\'owka, M. Lampa-Baczy\'nska, G. Malara, T. Szemberg, J. Szpond and H. Tutaj-Gasi\'nska,
\emph{Line arrangements with the maximal number of triple points},
Geom. Dedicata \textbf{180} (2016), 69--83; arXiv:1406.6662.

\bibitem{kuhne2024}
L. K\"uhne, T. Szemberg and H. Tutaj-Gasi\'nska,
\emph{Line arrangements with many triple points},
Rend. Circ. Mat. Palermo (2) \textbf{73} (2024), 2501--2512,
doi:10.1007/s12215-024-01047-x; arXiv:2401.14766.

\bibitem{balintova1994}
A. B\'alintov\'a and V. B\'alint,
\emph{On the number of circles determined by $n$ points in the Euclidean plane},
Acta Math. Hungar. \textbf{63} (1994), 283--289.

\bibitem{zhang2011}
R. Zhang,
\emph{On the number of ordinary circles determined by $n$ points},
Discrete Comput. Geom. \textbf{46} (2011), 205--211,
doi:10.1007/s00454-010-9286-8.

\bibitem{verificationarchive}
Denis Paliy,
\emph{Erdos506: exact verification archive for the minimum-circle problem},
Version 1.0.0, self-contained verification archive accompanying this preprint (2026).

\end{thebibliography}

